\section{Signed qualification and nonlinear transport: proofs}\label{supp:signed91}
\subsection{Proof of Theorem~\ref{thm:signed91}}
Fix $\theta$ and omit it from the notation. Because $b^-_t$ is valid for every original feasible action, taking infima gives
\[
 T_tf_{t+1}-f_t\ge\widehat T_t-f_t+b^-_t\ge\ell_t+b^-_t.
\]
For the actual feasible action, the numerical greediness and actor-specific upper bounds give
\[
 T_t^\pi f_{t+1}-f_t
 \le\widehat Q_t(\cdot,\pi_t)-f_t+b^{\pi,+}_t
 \le u_t+\gamma_t+b^{\pi,+}_t.
\]
Define real constants backwards by
\[
 a_T=d^-,\quad c_T=d^+,\qquad
 a_t=\ell_t+b^-_t+\beta_ta_{t+1},\quad
 c_t=u_t+\gamma_t+b^{\pi,+}_t+\beta_tc_{t+1}.
\]
At the terminal date $f_T+a_T\le g\le f_T+c_T$. If $V_{t+1}^*\ge f_{t+1}+a_{t+1}$, monotonicity and cash invariance, including under the infimum, imply
\[
 V_t^*=T_tV_{t+1}^*
 \ge T_tf_{t+1}+\beta_ta_{t+1}\ge f_t+a_t.
\]
Similarly, $V_{t+1}^\pi\le f_{t+1}+c_{t+1}$ implies
\[
 V_t^\pi=T_t^\pi V_{t+1}^\pi
 \le T_t^\pi f_{t+1}+\beta_tc_{t+1}\le f_t+c_t.
\]
Induction proves both comparisons. Subtracting at date zero gives $V_0^\pi-V_0^*\le c_0-a_0$, which expands to \eqref{eq:signed-loss91}. Feasibility and optimality give its lower bound zero. Infima need not be attained: the order comparison holds directly under infima, provided the stipulated original value recursion and feasible policy are defined. Every parameter is treated by the same argument; a supremum of valid bounds provides the stated classwise conclusion. No union over parameter vectors and no reference-policy occupation law is introduced.

For shifts $f'_t=f_t+k_t$, set $\widehat Q'_t=\widehat Q_t+\beta_t k_{t+1}$. Both numerical residual endpoints shift by $\beta_t k_{t+1}-k_t$; their difference is unchanged. Numerical action differences and original-minus-numerical action values are unchanged by cash invariance. Both terminal endpoints shift by $-k_T$. Thus every term in \eqref{eq:signed-loss91} is invariant. If the action-value discrepancy itself is a constant $k_t$, choose $b^-_t=b^{\pi,+}_t=k_t$ and its charge vanishes.

For the stationary assertion, the Bellman and fixed-policy maps are sup-norm contractions with modulus $\beta$: monotonicity and cash invariance imply sup-norm nonexpansiveness of each one-step evaluation. The finite-horizon argument with a bounded terminal comparison then has a terminal remainder of order $\beta^T$. Letting $T$ tend to infinity proves the displayed stationary bound. The asserted constant-shift invariance does not justify shifting an original preference primitive or an action-dependent immediate cost.

\subsection{Sharpness and the distinction between constant and statewise bias}
For $h>0$, take one decision date, terminal cost $g(y)=y$, discount one and zero immediate costs. The original successors of actions zero and one are respectively $h$ and zero; the numerical successors are both zero. Store $f_0=0$ and $f_1=g$, and let the numerical tie convention select action zero. All numerical residual and greediness terms vanish. The original-minus-numerical action values lie in $[0,h]$, the selected-action upper endpoint is $h$, and the implemented policy loses exactly $h$. Thus the signed charge $h$ is attained. The symmetric allowance $2h$ remains valid but is strictly larger. The example works for any monotone cash-invariant normalized risk evaluation because its successors are deterministic. Both laws can be singular, so no total-variation convergence is being used.

A statewise common-action bias is different. Let a root action choose state $L$ or $R$ at the next decision date. At that date there is only one feasible action, with original cost $2h$ in $L$ and zero in $R$. Let every numerical cost, stored continuation and terminal cost be zero, with discount one. Break the root numerical tie toward $L$. All numerical residual spans and greediness gaps vanish. At the second date the discrepancy is common to every action at its state, but equals $2h$ or zero depending on that state. The implemented policy loses $2h$. Cancelling the discrepancy merely because it is action independent at each state would falsely report zero loss. The signed uniform range correctly charges $2h$. This is precisely the earlier action's effect on future states; fixing the reference occupation would also lose it.

\subsection{Proof of Proposition~\ref{prop:risk-transport91}}
Let $(Y,\widehat Y)$ have any coupling of the two successor laws. For every real threshold $z$, the map $r\mapsto(r-z)_+$ is one-Lipschitz. Hence
\[
 \left|P(f-z)_+-\widehat P(f-z)_+\right|
 \le \E|f(Y)-f(\widehat Y)|\le K\E d(Y,\widehat Y).
\]
The same bound, divided by $1-\alpha$, is uniform in $z$ for the two functions minimized in \eqref{eq:risk-definitions91}. Infima of uniformly close functions are close. Taking the infimum over couplings proves the first bound, with approximating couplings sufficient when an optimum is not attained. At $\alpha=0$ the variational expression equals expectation: the expression is at least $Pf$ and equality holds for any threshold no larger than the bounded loss's infimum.

To verify that the tail factor cannot be uniformly reduced, take $f(y)=y$ on $[0,1]$, $\widehat P=\delta_0$ and $P=(1-h)\delta_0+h\delta_1$, with $0<h\le1$ and $\alpha=1-h$. Then $K=1$, $W_1(P,\widehat P)=h$, $\rho_P^\alpha(f)=1$ and $\rho_{\widehat P}^\alpha(f)=0$. The bound $Ke/(1-\alpha)=1$ is attained. In particular, a bound uniform up to risk level one cannot be deduced from Wasserstein distance alone with a fixed coefficient.

For the entropic assertion, choose $m\le f\le M$. For positive $\lambda$, the mean value theorem gives
\[
 |\exp(\lambda f(y))-\exp(\lambda f(\widehat y))|
 \le\lambda\exp(\lambda M)K d(y,\widehat y).
\]
Both exponential integrals are at least $\exp(\lambda m)$, and
$|\log s-\log t|\le |s-t|/\min(s,t)$ for positive $s,t$. Combining these facts, dividing by $\lambda$ and taking the infimum over couplings yields $K\exp\{\lambda(M-m)\}e$. Letting $m$ and $M$ approach the continuation's infimum and supremum proves \eqref{eq:risk-moduli91}. Translation invariance follows from the oscillation rather than the absolute level. The expectation limit follows from boundedness and expansion of the exponential integral; its displayed coefficient tends to $K$.

If $f(Y)\ge f(\widehat Y)$ almost surely, every positive-part integral is weakly larger and so is its minimized variational expression. Exponentiation and logarithm preserve the same order for the entropic evaluation. This gives the signed lower endpoint zero. If $|\widehat Q-c-\beta\rho_{\widehat P}(f)|\le\xi$, subtracting the numerical approximation places the action-value discrepancy in $[-\xi,\beta\kappa e+\xi]$ under that ordering, or in $[-\beta\kappa e-\xi,\beta\kappa e+\xi]$ without it. These are bounds for the original risk primitive; no statement is made about replacing one primitive by another.

\subsection{Nested objectives and the dated resource account}
Nested one-step evaluations and a static terminal evaluation differ even without decisions. With four equiprobable terminal losses $0,0,0,4$, group the first two and last two under two equiprobable intermediate states. At risk level one half, the intermediate evaluations are zero and four, and the nested value is four. Static average value at risk of the four terminal outcomes is two. Consequently a static objective needs its own state augmentation or a different analysis; applying the preceding nested recursion without that step changes the objective.

For an originally nested criterion, the preceding induction propagates a state-independent continuation allowance $a$ as $\beta_t a$. The local Wasserstein-to-risk coefficient $\kappa_t$ enters only the current discrepancy. Substituting the symmetric signed interval into \eqref{eq:signed-loss91} gives \eqref{eq:risk-budget91}; monotone-coupling information replaces the factor two on that transition component by one. This establishes the budget used for dated allocation. It does not bound the cost of obtaining $K_{t+1}$ or the number of labels, innovations or feasible-action searches. In particular, a growth of the verified continuation modulus with horizon is retained in $\kappa_t$ and cannot be removed by this algebra.

For a fixed-opponent best-response problem the identical induction compares a player's implemented complete strategy with the original best-response optimum. Applying it separately in every declared subgame gives the corresponding unilateral-deviation certificate. Each temporal self likewise retains its own original criterion and fixed future strategies. Monotone but non-cash-invariant aggregators do not satisfy the constant-shift step above and must instead use their certified absolute-error recursion. This establishes the advertised scope without conflating expectation, nested risk, static risk and strategic evaluation.
